\songtitle{Espejo Luminoso}

\tag*{2026}\tag{salsa}\tag{spanish-lyrics}

\begin{notes}
\end{notes}

\begin{style}
salsa, montuno, Spanish-English storytelling, male tenor lead, choir call and response, piano montuno riffs, conga tumbao, timbales cascara, bass tumbao, trumpet unison hits, tres montuno flourishes, horn stabs, live room bleed, 1970s analog warmth, abrupt dropouts, sparse verse to dense chorus, midtempo groove, bittersweet resolve, domestic tension
\end{style}

\begin{lyrics}

\songpart{Intro}
\annotation{Sets the stage, brass section playing a catchy, syncopated melody}

\songpart{Verse 1}
\annotation{Lead, male tenor vocals}\\
Hay quien busca en su morada\\
el refugio del descanso,\\
y hay quien vive su vivienda\\
como escudo y como manto.\\
Uno quiere solamente\\
silencio, calma, reposo;\\
el otro quiere una forma,\\
un espejo luminoso.

\songpart{Chorus}
\annotation{Choir} Hay dos tipos de personas\\
\annotation{Lead} uno vive en su refugio,\\
el otro vive su casa\\
como rostro y como pulso.\\
\annotation{Choir} Uno descansa sin culpa,\\
el otro carga el trabajo.

\songpart{Verse 2}
\annotation{Lead, male tenor vocals}\\
Un plato sobre la mesa,\\
una chaqueta en el suelo:\\
¿qué hay aquí que ver o corregir?\\
No entiendo ese desvelo.\\
Levantarme a recoger\\
sí arruina la vivencia;\\
mi casa es para vivirla,\\
no pa' perder la paciencia.

\songpart{Chorus}
\annotation{Choir} Hay dos tipos de personas\\
\annotation{Lead} uno vive en su refugio,\\
el otro vive su casa\\
como rostro y como pulso.\\
\annotation{Choir} Uno descansa sin culpa,\\
el otro carga el trabajo.

\songpart{Verse 3}
\annotation{Lead, male tenor vocals}\\
El plato sin lavar duele,\\
la chaqueta en el suelo pesa;\\
mi casa soy yo misma,\\
el desorden la atraviesa.\\
Y entonces toca elegir:\\
pelear o doblar el turno,\\
corregir lo que otros rompen,\\
callar y seguir el curso.

\songpart{Chorus}
\annotation{Choir} Hay dos tipos de personas\\
\annotation{Lead} uno vive en su refugio,\\
el otro vive su casa\\
como rostro y como pulso.\\
\annotation{Choir} Uno descansa sin culpa,\\
el otro carga el trabajo.

\songpart{Montuno}
\annotation{Choir} (¿Por qué el uno no arregla?)\\
\annotation{Lead} La casa es para vivirla\\
\annotation{Choir} (¿Por que no la arregla?)\\
\annotation{Lead} Porque arruina la vivencia\\
No hay que perder la paciencia

\annotation{Choir} (¿Por qué el otro se enoja?)\\
\annotation{Lead} Por corregir lo que otros rompen\\
\annotation{Choir} (¿Por qué no se sienta y descanza?)\\
\annotation{Lead} mi casa soy yo misma,\\
el desorden la atraviesa.

\annotation{Mambo, piano montuno riffs, conga tumbao,  bass tumbao, trumpet unison hits}

\songpart{Verse 4}
\annotation{Lead, male tenor vocals}\\
Uno ve un espacio libre,\\
el otro ve una herida;\\
no es maldad ni pereza,\\
es que no comparten vida.\\
El que descansa no finge,\\
el que carga no inventa;\\
habitan el mismo cuarto\\
en dos mundos sin encuentro

\songpart{Chorus - Outro}
\annotation{Choir} Hay dos tipos de personas\\
\annotation{Lead} uno vive en su refugio,\\
el otro vive su casa\\
como rostro y como pulso.\\
\annotation{Choir} Uno descansa sin culpa,\\
\annotation{Lead} el otro carga el trabajo.\\
\annotation{Choir} Hay dos tipos de personas\\
Uno descansa sin culpa,\\
el otro carga el trabajo.\\
el refugio del descanso,\\
y hay quien vive su vivienda\\
como escudo y como manto.

\end{lyrics}

